\chapter{Méthodes d’alignement}

Les premières méthodes d’alignement automatique de phrases reposaient principalement sur
des approches statistiques basées sur la longueur des segments. L’algorithme proposé par
\citet{gale_1991} constitue un exemple fondateur. Il exploite la corrélation entre la
longueur des phrases dans deux langues pour établir des correspondances probabilistes.
Ces méthodes ont posé les bases des approches modernes, aujourd’hui enrichies par des
représentations vectorielles et des modèles neuronaux.

L'une des directions clés de la recherche pour le TAL moderne consiste à apprendre dans un environnement se basant sur les langues peu dotées.
\citep{hedderich_2021, artemova_2023}. Cette affirmation vient confirmer ce que nous avons présenté
auparavant dans les chapitres de l'état de l'art. Maintenant la question n'est plus d'expliquer les
problématiques et définitions de notre environnement, mais plutôt les solutions élaborées par le passé.
Ces méthodes d'alignement sont très pratiques et il est préférable que nous les présentions par projet.

Il est important de noter que les données parallèles pour les langues peu dotées sont souvent en petit nombre, la qualité
de ces données est donc primordial. Il faut donc limiter au maximum les données présentant des phrases atypiques,
ayant des rapport de longueur trop grand ou trop petit, ou un texte se chevauchant entre langue source et langue cible.
Ces phrases risquent de se retrouver dans les données parallèles et posent un problème dans le contexte de la traduction
automatique par réseau neuronal moderne, qui dépend énormément de la bonne qualité des ressources
d'entraînement \citep{tan_2022}. Des méthodes non-supervisées ont été proposées pour pallier ce problème, comme
l'entraînement sur des corpus monolingues seuls \citep{artetxe_2018, lample_2018}. Cependant,
\citet{steingrimsson_2021} nous expliquent que d'autres travaux par \citet{kim_2020} montrent que des méthodes
supervisées et semi-supervisées avec de petits corpus parallèles seront toujours meilleures que des méthodes
non-supervisées pour beaucoup de langues. Le manque de ressources parallèles existantes pour les langues peu
dotées implique qu'il faudrait plutôt utiliser des corpus comparables en supplément de corpus parallèles.
Ces données permettraient l'entraînement de réseaux neuronaux modernes afin d'atteindre une meilleure qualité sur des modèles plus récents
\citep{steingrimsson_2023}.

\section{Principes fondamentaux}

Les réseaux de neurones récurrents (RNN) et leurs variantes à mémoire à long terme (LSTM) ont constitué
une étape importante dans le traitement des séquences textuelles avant l'apparition des \gls{transformers}. Ces architectures
ont été largement utilisées dans les premières méthodes d'alignement basées sur des réseaux neuronaux, notamment
pour modéliser les dépendances entre phrases source et cible. Les méthodes d'alignement neuronal ont progressivement
remplacé les approches purement statistiques en permettant de capturer des similarités sémantiques au-delà des
simple heuristiques de longueur ou de fréquence de mots. Cette évolution a ouvert la voie aux approches modernes
fondées sur les transformers et les plongements multilingues, qui constituent le socle des méthodes d'alignement
actuelles.

\subsection{RNN et LSTM}

Les RNN (\textit{Recurrent Neural Networks}) sont des réseaux neuronaux conçus pour traiter des
données séquentielles, comme du texte ou des paroles.
Contrairement aux réseaux neuronaux classiques (feed-forward), les RNN possèdent une mémoire interne.
L’état calculé à un instant donné dépend non seulement de l’entrée actuelle, mais aussi de l’état précédent.
\[ h_{t}=f(W_{x_{t}}+U_{h_{t-1}}+b) \]
où:
\begin{itemize}
    \item \( x_{t} \) est l'entrée au temps \( t \)
    \item \( h_{t-1} \) est l'état précédent
    \item \( W, U, b \) sont des paramètres appris
\end{itemize}
L'avantage principal des RNN est de modéliser naturellement l’ordre des mots.
Cependant ils présentent deux limites importantes{ :}
\begin{itemize}
    \item Le modèle a du mal à apprendre des dépendances longues
    \item Chaque état dépendant du précédent, l'entraînement est fortement ralenti
\end{itemize}

Les LSTM (\textit{Long Short-Term Memory}) ont été introduits par \citet{hochreiter_1997} pour résoudre le problème des dépendances longues.
Au lieu d’avoir une seule mémoire simple, les LSTM introduisent une mémoire cellulaire contrôlée par des portes.
Les portes décident quelles informations sont à garder et à transmettre.
On peut voir les LSTM comme une mémoire intelligente qui retient les informations importantes et oublie les informations
qui ne sont plus pertinentes, ce qui permet une mise à jour du contexte contrôlé.
Malgré ces améliorations, les LSTM restent lents et deviennent lourds dans le cas de séquences longues.

\subsection{Principe de transformer}

Le transformer est une architecture de réseau neuronal introduite par \citet{vaswani_2017}.
Cette architecture a profondément transformé le traitement automatique des langues en remplaçant
les architectures récurrentes (RNN, LSTM) par un mécanisme reposant entièrement sur l’attention.
Contrairement aux modèles séquentiels classiques, le transformer ne traite pas les mots d’une phrase
un par un. Il traite l’ensemble de la séquence simultanément. Son idée centrale est la suivante{ :}
pour comprendre un mot dans une phrase, il faut pouvoir mesurer son importance relative par
rapport à tous les autres mots de la phrase. Cette opération est réalisée via un mécanisme appelé auto-attention.

Chaque mot est transformé en un vecteur numérique (\textit{embedding} ou plongement en français).
À partir de ce plongement, trois vecteurs sont calculés{ :}
\begin{itemize}
    \item Query (Q)
    \item Key (K)
    \item Value (V)
\end{itemize}
À partir de ces vecteurs nous calculons l'attention, définie de cette manière{ :}
\[ Attention(Q,K,V)=softmax(\frac{QK^{T}}{\sqrt{d_{k}}})V \]
où \( d_{k} \) est la dimension des vecteurs Key (facteur de stabilisation).
Ce calcul utilise la fonction \textit{\gls{fonction softmax}}, une fonction
mathématique utilisée pour transformer un vecteur de valeurs réelles (appelées logits)
en une distribution de probabilités.

Le transformer n'utilise pas qu'une seule attention mais plusieurs "têtes" d'attention. Qui
permettent de capter plusieurs types de relations (comme les relations syntaxiques ou sémantiques).
Ces sorties sont ensuite concaténées puis projetées dans un nouvel espace vectoriel.
Le transformer traitant tous les mots en même temps, il ajoute aux plongements des vecteurs
de position. Ceux-ci permettent de savoir le positionnement des mots les uns par rapport aux autres 
et la distance entre ceux-ci.

Ce modèle présente plusieurs avantages, comme{ :}
\begin{itemize}
    \item Une meilleure gestion des dépendances longues
    \item Une vitesse d'entraînement bien plus rapide sur GPU due à la parallélisation des calculs
    \item La représentation contextuelle des mots
\end{itemize}
     
Cependant il est à noter que les transformers présentent aussi certains désavantages, comme{ :}
\begin{itemize}
    \item Un besoin de données en grandes quantités (complexe, dans notre cas par exemple)
    \item Une architecture beaucoup plus lourde tournant difficilement sur de plus petites machines.
\end{itemize}

Ces principes fondamentaux sont la base de la majorité des projets de TAL moderne. De nos jours, la vaste majorité
des projets se basent sur le principe de transformer, celui-ci étant l'un des plus performants pour ce domaine
d'application. Les projets sur lesquels nous allons nous baser dans ce mémoire se basent sur les transformer de
manière à être le plus efficace dans les tâches de traitement linguistique.

\section{BERT et SBERT}

BERT \citep{devlin_2019} (\textit{Bidirectional Encoder Representations from Transformers}) et RoBERTa
\citep{liu_2019} sont des modèles de transformer encodeur pré-entraînés. Un transformer encodeur est un transformer permettant
de traiter des séquences d'entrée et de générer des représentations viables dans un contexte donné. Les décodeurs, quant à eux
vont générer une sortie en se basant sur les représentations encodées. Les décodeurs s'assurent que chaque token généré
se base sur les tokens précédents. À l'époque de leur création, les modèles de type BERT
étaient très performants pour les tâches de comparaison de paires de mots telles que la \gls{similarité sémantique}.
En français, deux modèles principaux se basent sur BERT{ :} CamemBERT 
\citep{martin_2020} et FlauBERT \citep{le_2020}. Ces deux modèles sont entraînés sur des données d'un corpus
français hétérogène, BERT ne se basant uniquement que sur des données d'entraînement en anglais.

BERT est entraîné de manière auto-supervisée en utilisant 2 tâches principales{ :}
\begin{itemize}
    \item MLM (\textit{Masked Language Modeling}){ :} on masque aléatoirement environ 15 \% des mots d’une phrase,
    avant de demander au modèle de prédire les mots masqués en utilisant le contexte entourant le mot caché.
    \item NSP {\textit{Next Sentence Prediction}}{ :} on donne au modèle deux segments de texte, puis on demande de
    prédire si le second segment suit réellement le premier dans le texte d’origine.
\end{itemize}
Cet entraînement génère des représentations vectorielles pour chaque token et des tokens de classification (CLS-Token)
donnant une représentation globale de la phrase.

Plus tard \citet{reimers_2019} présentent SBERT ou sentence-BERT qui se base sur le modèle pré-entraîné 
BERT. L'objectif de SBERT est de transformer BERT en un modèle capable de générer directement des plongements
de phrases de taille fixe, tout en préservant leur signification sémantique. Plutôt que de traiter
chaque phrase de manière indépendante et de concaténer les sorties, SBERT utilise un réseau siamois
qui encode les deux phrases simultanément, permettant un calcul efficace de similarité.

Pour ce faire, SBERT va d'abord prendre le résultat de sortie de BERT afin d'en dériver un plongement de phrases de taille
fixe. On fait ensuite la médiane des vecteurs de sorties des 2 phrases, établis en utilisant la méthode de calcul
CLS-Token de BERT. On utilise ensuite des structures de \gls{réseau siamois} et triplets
dérivant des \gls{plongement lexicaux}, technique représentant un mot par une liste de nombres
qui encode les mots par similarité d'utilisation, ensuite comparé par \gls{similarité cosinus}.
Ces plongements sont ensuite concaténés et classifiés en utilisant la fonction softmax.
Celle-ci permet d'optimiser l'entropie croisée. Cette amélioration réduit le temps pour trouver les paires de phrases les plus
similaires à 5 secondes au lieu de {65 heures} tout en gardant la précision de BERT. La figure \ref{SBERT_schema} présente le fonctionnement
d'un réseau siamois SBERT.

Dans ce mémoire, l'utilisation de SBERT s'impose pour deux raisons principales. D'une part, BERT
nécessite de comparer chaque paire de phrases de manière individuelle, ce qui rend la recherche
dans un large corpus prohibitif en temps de calcul. D'autre part, le token [CLS] de BERT n'a pas
été explicitement optimisé pour la similarité de phrases, contrairement à l'approche de SBERT qui
a été entraînée spécifiquement sur des tâches de similarité sémantique.
De plus, d'après \citet{reimers_2019} SBERT surpasse les autres méthodes de \gls{plongement de phrases} (sentence embedding)
pour des tâches de similarité sémantique communes comme l'identification de paraphrase, la recherche sémantique
ou encore l'entailment/contradiction (détermine si une phrase A implique une phrase B, contredit B
ou n'ont pas de lien).

\begin{figure}[H]
    \centering
    \includegraphics[width=\textwidth]{SBERT_schema.png}
    \\
    \smallskip
    \small{Schéma se basant sur la figure 2 du travail d'\citet{ortakci_2024} et sur la figure 6 du travail de \citet{mokoatle_2022}}
    \smallskip
    \caption{Schéma du fonctionnement de SBERT (réseau siamois).}
\label{SBERT_schema}
\end{figure}

\newpage

\section{Paracrawl et bitextes}

\href{https://paracrawl.eu/}{ParaCrawl} \citep{banon_2020} est un robot d'exploration du web
(et outil open-source) permettant le développement de larges corpus parallèles de textes traduits
pour les langues officielles de l'Union européenne. Des outils tels que Hunalign \citep{varga_2007} sont
utilisés pour l'alignement de phrases pour des paires de documents dans le projet ParaCrawl
\citep{tan_2022}. Hunalign est un projet qui utilise des {« }heuristiques telles que la longueur des phrases
ou la fréquence des mots pour trouver des paires de phrases dans deux documents.{ »}
\footnote{{« }[...] heuristics such as sentence length and word frequency to find matching sentence pairs in
two documents.{ »} \citep{tan_2022} notre traduction}
\citep{tan_2022}

ParaCrawl mine des bitextes en passant par 4 étapes majeures \citep{tan_2022} { :}
\begin{enumerate}
\item Exploration de sites web{ :} pour collecter du contenu provenant du web.
\item Alignement documentaire{ :} trouve les documents qui s'alignent dans
différentes langues.
\item Alignement des phrases{ :} trouver des paires de phrases correspondantes
alignées dans les documents alignés précédemment. Alignement au niveau du contenu.
\item Filtrage des phrases{ :} filtrer les paires de phrases de mauvaise qualité et
utiliser le reste comme des données parallèles propres.
\end{enumerate}

Le projet de \citet{tan_2022} s'intéresse tout particulièrement aux étapes 3 et 4 et vise à les
améliorer par l'utilisation de meilleures fonctions d'évaluation de qualité d'alignement. Ils essaient 
principalement de réaliser
cette amélioration pour une paire de langues peu dotées (pachto et khmer). Ainsi \citet{tan_2022}
décident d'affiner le fonctionnement d'un modèle de transformer de phrases 
afin d'aligner et de filtrer
des phrases pour les langues peu dotées recherchées.
\citet{tan_2022} utilisent un système d'apprentissage
auto-supervisé \citep{chen_2020, henderson_2017} afin d'affiner le fonctionnement des modèles BERT
\citep{devlin_2019} et \allowbreak SBERT \citep{reimers_2019}.
L'objectif est de rendre le modèle plus précis sur une tâche particulière en le ré-entraînant sur un jeu de données
plus petit. Pour affiner le fonctionnement de ces transformers,
\citet{tan_2022} décident de construire des paires positives et négatives pour entraîner les modèles par
\gls{MNRL} (\textit{Multiple Negative Ranking Loss}).
Dans cette approche, chaque phrase source est associée à sa vraie traduction (positive).Toutes les autres phrases
cibles du batch sont considérées comme des négatives pour cette phrase source. Le modèle est entraîné à maximiser
la similarité de la vraie paire tout en minimisant la similarité avec toutes les autres paires du batch, ce qui
permet d'exploiter la diversité des comparaisons sans nécessiter de négatives explicites.

Le but est de faire en sorte que le modèle maximise la différence entre ces paires positives et négatives afin
de faciliter sa reconnaissance de phrases similaires tout en écartant les phrases inutilisables.
Les résultats de \citet{tan_2022} montrent une meilleure production de données d'entraînement (ici des bitextes)
pour la traduction par réseau neuronal qui aura lieu en aval de la récupération des données d'entraînement.
Ces données sont récupérées grâce au modèle SBERT \citep{reimers_2019}
combiné avec Vecalign \citep{thompson_2019} pour aligner les phrases présentes dans les données puis les filtrer à l'aide d'un score basé sur la marge. 
Vecalign est une méthode d'alignement de phrases introduite par \citet{thompson_2019} qui se distingue par sa complexité
linéaire en temps et en espace par rapport au nombre de phrases alignées. Contrairement aux méthodes précédentes qui nécessitaient un système
de traduction automatique, Vecalign ne requiert que des plongements de phrases bilingues pour établir des alignements. Cette approche
permet d'aligner des corpus à grande échelle avec une efficacité supérieure aux méthodes passées, tout en nécessitant moins de
ressources computationnelles.
Le score réalisé avec Vecalign et SBERT correspond à la différence entre la similarité (ou distance) d’une paire de phrases considérée et la similarité
entre cette même phrase et la phrase la plus proche parmi les négatives. Cette dernière permet de pallier les problèmes
d'inconsistances d'échelles donnant un résultat plus fiable que d'autres méthodes.

\newpage

\section{CLIR et plongement lexical}

Le \gls{CLIR} (\textit{crosslingual information retrieval} - 
récupération d'information inter-langue) désigne l’ensemble des méthodes permettant
d’effectuer une recherche d’information dans une langue différente de celle de la requête.
Un utilisateur formule une requête dans une langue source, et le système
identifie des documents pertinents rédigés dans une langue cible. Les approches de CLIR
reposent sur la traduction de la requête ou des documents via des dictionnaires
bilingues ou des systèmes de traduction automatique. Plus récemment, elles utilisent des
représentations vectorielles multilingues qui projettent différentes langues dans un espace sémantique partagé,
permettant ainsi de comparer directement des textes. Le CLIR
constitue une base théorique importante pour l’alignement multilingue, puisqu’il vise déjà
à mesurer la similarité sémantique entre des textes dans des langues différentes.

Dans le cadre de ce mémoire, utiliser
un aligneur de documents constitue une étape stratégique. Avant de rechercher
des correspondances au niveau des phrases ou des segments, il est nécessaire de restreindre
l’espace de recherche à des documents traitant du même sujet. L’alignement documentaire permet
ainsi de regrouper des textes susceptibles de contenir des segments parallèles, ce qui réduit
considérablement le bruit et le coût de calcul.

FaDA (\textit{Fast Document Aligner} - aligneur de document rapide) \citep{lohar_2016} est un outil d'alignement
de documents multilingues se basant sur un algorithme CLIR. Un
aligneur de documents inter-langue vise à extraire des candidats probables de documents
alignés venant d'un corpus comparable multilingue \citep{lohar_2016}. D'autres projets présentent des méthodes
se basant sur différents niveaux (titre de document, mots ou encore caractères dans certaines langues) comme les travaux de
\citet{yang_2003} ou encore le système de requête sur bases de données pour des lexiques bilingues présentés
par \citet{munteanu_2005}.

\citet{lohar_2016} proposent une méthode de plongement lexical vectoriel.
Plus des mots sont utilisables dans le même contexte, plus leurs vecteurs sont proches. Cette stratégie permet
de vérifier la similarité des mots d'une manière optimisée.
À cette méthode, \citet{lohar_2016} incorporent une approche basée sur la
CLIR pour estimer la similarité sémantique entre deux documents de langue source et langue cible. Leur système
commence par indexer les documents de langue source et langue cible. Les documents sources sont utilisés pour
faire une pseudo-requête composée d'une fraction des termes constituant les documents dans la langue source.
Les termes choisis sont trouvés à partir d'un score basé sur une équation où {« }$tf(t,d)$ dénote la fréquence
d'un terme $t$ dans un document $d$, $len(d)$ dénote la longueur de $d$ et $N$ et $df(t)$ dénotent respectivement
le nombre total de documents et le nombre de documents dans lesquels $t$ apparaît. De plus, $\tau(t,d)$
représente le score de sélection des termes.{ »} \citep{lohar_2016} \footnote{{« }$tf(t, d)$ denotes the term frequency of a
term $t$ in document $d$, $len(d)$ denotes the length of $d$, and $N$ and $df(t)$ denote the total number of documents and the
number of documents in which $t$ occurs, respectively. Furthermore, $\tau(t, d)$ represents the term-selection score{ »} Notre traduction}.

\[ \tau(t,d)=\lambda\frac{tf(t,d)}{len(d)}+(1-\lambda)log(\frac{N}{df(t)}) \]

Les termes en langue source sont ensuite traduits via un dictionnaire vers la langue cible indexée avant
d'être utilisés pour réaliser une comparaison entre la liste de termes et les documents dans la langue cible.
Cette méthode de comparaison permet d'extraire les meilleurs documents et de sélectionner celui avec le
meilleur score pour avoir le meilleur candidat pour l'alignement.

Plus précisément, \citet{steingrimsson_2021} utilisent FaDA ainsi qu'une technique de scoring de phrases candidates basée sur
le plongement contextualisé et d'autres alignements par mot de haute précision. Enfin, un classifieur
sélectionne des paires de phrases en se basant sur ces scores. 

\citet{steingrimsson_2021} utilisent FaDA pour produire un ensemble de candidats
(typiquement 10 phrases cibles par phrase source, dans les deux directions de traduction). Pour filtrer ces candidats,
ils combinent deux scores complémentaires. D'une part un score basé sur les plongements de phrases LaBSE \citep{feng__2022},
un modèle BERT considéré comme meilleur que LASER \citep{artetxe_2019} (autre modèle produisant des
embeddings multilingues comparables directement par similarité cosinus), celui-ci peut être utilisé pour le
calcul de similarité entre des phrases. Il se base sur une fonction de perte de type \textit{translation ranking} et un
bi-\textit{encoder}. La \textit{translation ranking loss} consiste à présenter au modèle une phrase dans la langue source accompagnée de sa traduction correcte
(positif) et de plusieurs traductions incorrectes tirées d'autres documents (négatifs). Le modèle est entraîné à maximiser la similarité
de la paire correcte tout en minimisant celle des paires incorrectes. L'architecture bi-encoder encode la phrase source et
la phrase cible indépendamment dans deux encodeurs identiques, puis calcule leur similarité par produit scalaire. Cette approche est plus efficace que
l'encodeur-décodeur classique car les plongements peuvent être précalculés et indexés pour une recherche rapide.
D'autre part \citet{steingrimsson_2021} utilisent un score original nommé WAScore, calculé à partir d'alignements de mots de haute précision obtenus via
CombAlign \citep{steingrimsson_combalign_2021}. Celle-ci combine trois outils d'alignement (AWESoME \citep{dou_2021},
eflomal \citep{ostling_2016}, fast\_align \citep{dyer_2013}) en visant la précision plutôt que le rappel.
Le WAScore est défini comme le produit des proportions de mots alignés dans les deux directions (\( (sa/s) \times (ta/t) \)), où \(s\) et \(t\)
sont les longueurs des phrases source et cible, et \(sa\) et \(ta\) le nombre de mots alignés. Cette formulation pénalise les paires de longueurs
très disparates et favorise les phrases dont la plupart des mots sont mutuellement alignés. Ces deux scores (LaBSE \citep{feng__2022} et WAScore) sont ensuite
combinés par un classifieur de régression logistique entraîné sur des paires étiquetées manuellement, qui prédit si une paire de phrases est
parallèle ou non. Le caractère utile des phrases parallèles extraites est confirmé par l'entraînement de systèmes de traduction neuronale.
L'ajout des données extraites depuis Wikipedia augmente le score BLEU de 9 points par rapport au corpus
parallèle de référence seul.

\section{Réutilisation de segments parallèles}

\citet{steingrimsson_2023} ont réalisé des travaux de recherche sur la ré-examination de données
d'entraînement parallèles rejetées car elles ne passent pas un test d'équivalence sémantique (rapport de longueur impair,
texte se chevauchant entre langue source et langue cible, ...) en les segmentant avant d'extraire sémantiquement d'autres segments
bilingues équivalents. Plus précisément, \citet{steingrimsson_2023} appliquent un découpage basé sur la ponctuation forte (points,
points d'exclamation, points d'interrogation) pour diviser les phrases en segments plus courts. Ensuite,
ils comparent chaque segment avec son équivalent potentiel dans l'autre langue en utilisant LASER \citep{artetxe_2019}, LaBSE \citep{feng__2022} et WAScore.
Les segments qui atteignent un score de similarité suffisant sont conservés comme données parallèles exploitables,
même si la phrase complète avait été rejetée. La méthode se base sur l'entraînement de modèles de traduction
automatique neuronale (\textit{Neural Machine Translation} (NMT)) séparés pour les deux directions de la traduction.
Dans leur cas, de l'anglais au bengali et inversement. \citet{steingrimsson_2023} utilisent ces 3 méthodes de scoring
afin de donner un score à chacune des paires de phrases parallèles.

Les paires de phrases sont ensuite filtrées et classées avec un classifieur à \gls{régression logistique}
utilisant le score concaténé précédemment établi à l'aide de LASER, LaBSE et WAScore.
Une fois classées, un score \gls{BLEU} est établi sur les paires extraites d'un corpus anglais-bengali déjà réalisé par
\citet{ramesh_2022}. Le score BLEU établi est ici utiliser pour évaluer la qualité des traductions générées par divers modèles
créé par \citet{steingrimsson_2023} se basant sur le corpus Samanantar de \citet{ramesh_2022} entre elles.
Le score BLEU est une méthode de calcul de précision{ :}

Sous la forme{ :}
\[ BLEU=min(1,exp(1-\frac{reference-longueur}{sortie-longueur}))(\prod_{i=1}^{4}precision_{i}^{1/4}) \]
Où{ :}
\begin{itemize}
    \item $reference$ désigne la traduction de référence produite par un humain, considérée comme correcte ou acceptable.
    \item $sortie$ correspond au texte généré par le système de traduction automatique.
    \item $longueur$ correspond au nombre total de mots (ou tokens) présents dans un texte.
    \item Dans $precision_{i}^{1/4}$, $1/4$ correspond à la taille des n-grammes allant de 1 à 4
\end{itemize}
et $precision_{i}$ correspond à{ :}
\[ precision_{i}=\frac{\sum_{snt \in Cand-Corpus}\sum_{i \in snt}min(m^{i}_{cand},m^{i}_{ref})}{w^{i}_{t}=\sum_{snt' \in Cand-Corpus}\sum_{i' \in snt'}m^{i'}_{cand}} \]
Où{ :}
\begin{itemize}
    \item $m^{i}_{cand}$ est le nombre d'i-grammes correspondant à la traduction de référence.
    \item $m^{i}_{ref}$ est le nombre d'i-grammes dans la traduction de référence.
    \item $w^{i}_{t}$ est le nombre d'i-grammes dans la traduction automatique.
    \item $snt$ signifie \textit{sentence} soit{ :} $snt \in Cand-Corpus$ l'ensemble des phrases 
    produites par le système de traduction automatique.
\end{itemize}

Le score BLEU final est compris entre zéro et un représentant
l'exactitude d'une traduction. Le score est souvent représenté sous forme de pourcentage, au-dessus de 30 \%,
la traduction est considérée comme bonne. Il se base sur une correspondance de \gls{n-gramme} et une pénalité de concision.
Elle pénalise les traductions générées trop courtes par rapport
aux traductions de référence. La correspondance de n-gramme compte le nombre de n-grammes (allant de unigrammes à quadrigrammes)
correspondant à leurs équivalents dans la traduction de référence.

En utilisant cette méthode, \citet{steingrimsson_2023} ont améliorer la qualité du corpus d'entraînement tout 
en augmentant le score BLEU de 1,6 dans le cas de la traduction anglais-bengali. Le but étant de
démontrer qu'en ajoutant des données partiels supplémentaires à un corpus pouvait améliorer la qualité des traductions de celui-ci.
Ce système, entraîné sur moins de données, montre de meilleurs 
résultats que des tests se basant sur un modèle de transformer référentiel.
Par « modèle de transformer référentiel », \citet{steingrimsson_2023} désignent
un modèle transformer standard entraîné sur l'intégralité des données parallèles disponibles
(sans réutilisation des segments rejetés). Il sert de base pour comparer
les performances de leur approche optimisée.

Ce dernier est entraîné avec Fairseq \citep{ott_2019}. Fairseq est un outil open-source pour {« }l'entraînement de modèles de traduction, de résumé, 
de modélisation linguistique et d'autres tâches de génération. { »}\footnote{{« }train custom models for 
translation, summarization, language modeling, and other text generation tasks.{ »} \citep{ott_2019} 
notre traduction} \citep{ott_2019}. \citet{steingrimsson_2023} l'utilise ici car il s'agit de la bibliothèque de référence pour l'entraînement
de modèles de traduction automatique neuronale à grande échelle.
Elle offre une implémentation optimisée du transformer et a été utilisée dans la majorité des
travaux de recherche sur la traduction neuronale, ce qui permet de comparer directement les résultats avec la littérature existante.
Dans leur méthode, Fairseq utilise un vocabulaire de 32 000 termes encodés par paires de byte. 
Celui-ci était évalué automatiquement selon le score BLEU et les mêmes ensembles de données que les modèles 
principaux.

\vspace{5mm}

Cette revue des méthodes d'alignement, des fondamentaux statistiques aux approches neuronales modernes,
illustre comment la recherche a progressivement déplacé l'axe des heuristiques de longueur vers des
représentations sémantiques multilingues. La section suivante analyse l'état de l'art de manière critique
et retrace les limites, avant de présenter les données exploitées et la méthodologie que nous avons
développée pour identifier des fragments parallèles dans des corpus comparables de langues régionales de France.